\documentclass[11pt]{article}

\usepackage[margin=2.5cm]{geometry}
\usepackage{amsmath,amssymb,amsthm}
\usepackage{booktabs}
\usepackage{array}
\usepackage{hyperref}

\newcommand{\dd}{\mathrm{d}}
\newcommand{\cod}{\delta}
\newcommand{\tr}{\operatorname{tr}}
\newcommand{\hodge}{\star}
\newcommand{\Hspace}{\mathfrak{h}}
\newcommand{\R}{\mathbb{R}}
\newcommand{\curltwo}{\operatorname{curl_{2D}}}
\newcommand{\skewgrad}{\operatorname{skewgrad}}
\newcommand{\grad}{\operatorname{grad}}
\newcommand{\curl}{\operatorname{curl}}
\renewcommand{\div}{\operatorname{div}}
\newcommand{\code}[1]{\texttt{#1}}
\newcommand{\ip}[2]{\left(#1,\,#2\right)}

\newtheorem{proposition}{Proposition}
\theoremstyle{remark}
\newtheorem{remark}{Remark}

\title{The Hodge Laplacian in DIMSWE:\\ differential forms, vector proxies and sign conventions}
\author{DIMSWE developers}
\date{}

\begin{document}
\maketitle

\begin{abstract}
This note documents the elliptic Hodge Laplacian problems solved by \code{dimswe.elliptic}.
It states the continuous problem in the language of differential forms, the mixed weak
formulation of Arnold, Falk and Winther, and the vector-calculus proxies used for every
(dimension, complex, form degree) combination. It fixes every sign convention used in the code:
the exterior derivative and codifferential of each proxy, the boundary pairings, the proxies of
constant forms, the sign of the mixed system, and the treatment of boundary conditions, harmonic
forms and their two constraint mechanisms. Section~\ref{sec:code} maps each formula to the
function that implements it.
\end{abstract}

\tableofcontents

%=================================================================================================
\section{Differential forms on box-shaped domains}
%=================================================================================================

Let $\Omega\subset\R^n$, $n\in\{1,2,3\}$, be a box $\prod_{i=1}^n [a_i,b_i]$ in which each
coordinate direction is either \emph{bounded} or \emph{periodic} (identified ends). We write
$\Lambda^k(\Omega)$ for the smooth differential $k$-forms, $0\le k\le n$. With Cartesian
coordinates $x_1,\dots,x_n$ (in the code: \code{x}, \code{y}, \code{z}) every $k$-form is
\[
  \omega = \sum_{|I|=k} \omega_I \, \dd x^I, \qquad
  \dd x^I = \dd x_{i_1}\wedge\cdots\wedge\dd x_{i_k}, \quad I=\{i_1<\dots<i_k\}.
\]
\emph{Index sets are always taken in increasing order} $x<y<z$; this matters for the signs of
the constant-form proxies in Section~\ref{sec:constant}.

\paragraph{Exterior derivative.} $\dd:\Lambda^k\to\Lambda^{k+1}$, $\dd\circ\dd=0$, with
$\dd\Lambda^n=0$. The sequence
\[
  0 \to \Lambda^0 \xrightarrow{\dd} \Lambda^1 \xrightarrow{\dd} \cdots \xrightarrow{\dd} \Lambda^n \to 0
\]
is the de Rham complex.

\paragraph{Inner product and Hodge star.} The Euclidean metric gives the pointwise inner product
in which the $\dd x^I$ are orthonormal, the $L^2$ inner product $\ip{\omega}{\eta} =
\int_\Omega \langle\omega,\eta\rangle\,\dd x$, and the Hodge star $\hodge:\Lambda^k\to\Lambda^{n-k}$
with $\omega\wedge\hodge\eta = \langle\omega,\eta\rangle\,\dd x_1\wedge\cdots\wedge\dd x_n$.

\paragraph{Codifferential.} $\cod:\Lambda^{k}\to\Lambda^{k-1}$ is the formal $L^2$ adjoint of
$\dd$: for $w\in\Lambda^{k}$ and $z\in\Lambda^{k+1}$,
\begin{equation}\label{eq:ibp-forms}
  \ip{\dd w}{z} = \ip{w}{\cod z} + \int_{\partial\Omega} \tr w \wedge \tr \hodge z ,
\end{equation}
where $\tr$ is the pullback to the boundary. We never need an explicit formula for $\cod$ in
terms of $\hodge$; in the code $\cod$ is \emph{defined} through its proxy formulas
(Table~\ref{tab:operators}), which are verified to satisfy \eqref{eq:ibp-forms}.

\paragraph{Hodge Laplacian.} $\Delta_H = \dd\cod + \cod\dd : \Lambda^k\to\Lambda^k$
(non-negative; for $k=0$ it is $-\Delta$). The problem is: given $f\in\Lambda^k$, find
$u\in\Lambda^k$ with
\begin{equation}\label{eq:strong}
  (\dd\cod + \cod\dd)\,u = f - p \quad\text{in }\Omega, \qquad u\perp\Hspace^k, \qquad p\in\Hspace^k,
\end{equation}
plus boundary conditions (Section~\ref{sec:bc}). Here $\Hspace^k$ is the space of harmonic
$k$-forms (Section~\ref{sec:harmonic}); $p=P_{\Hspace}f$ is the harmonic part of the data that
the operator cannot reach. When $\Hspace^k=\{0\}$, $p$ is absent and \eqref{eq:strong} is the
usual Hodge Laplace equation.

%=================================================================================================
\section{Mixed weak formulation and the sign convention of the code}\label{sec:weak}
%=================================================================================================

Introduce $\sigma=\cod u\in\Lambda^{k-1}$ (absent for $k=0$). Testing $\sigma=\cod u$ with
$\tau\in\Lambda^{k-1}$ and using \eqref{eq:ibp-forms} with $(w,z)=(\tau,u)$ gives
$\ip{\sigma}{\tau} = \ip{u}{\dd\tau} - \int_{\partial\Omega}\tr\tau\wedge\tr\hodge u$. Testing
\eqref{eq:strong} with $v\in\Lambda^k$ and using \eqref{eq:ibp-forms} with $(w,z)=(v,\dd u)$
gives $\ip{\cod\dd u}{v} = \ip{\dd u}{\dd v} - \int_{\partial\Omega}\tr v\wedge\tr\hodge\dd u$.

\paragraph{The code's form.} The $\sigma$ equation is \emph{multiplied by $-1$}, so that the
system operator is symmetric (this is what makes the PETSc nullspace equal to the transpose
nullspace and allows MINRES):
\begin{equation}\label{eq:mixed}
\boxed{
\begin{aligned}
  -\ip{\sigma}{\tau} + \ip{u}{\dd\tau} &= \textstyle\int_{\Gamma_n} B_{k-1}(\tau, g_u)
     && \forall\tau\in V^{k-1},\\
  \ip{\dd\sigma}{v} + \ip{\dd u}{\dd v} + \ip{p}{v} &= \ip{f}{v} + \textstyle\int_{\Gamma_n} B_k(v, g_{\dd u})
     && \forall v\in V^{k},\\
  \ip{u}{q} &= 0 && \forall q\in\Hspace^k,
\end{aligned}}
\end{equation}
with $B_j(w,z)\,\dd s = \tr w\wedge\tr\hodge z$ the boundary pairing of \eqref{eq:ibp-forms} for
$\dd$ acting on $j$-forms (Table~\ref{tab:operators}). The terms on essential faces $\Gamma_e$
drop out because the test functions have vanishing trace there. Absent pieces:
$\sigma,\tau$ and the first row for $k=0$ (primal problem); $\ip{\dd u}{\dd v}$ and the
$B_k$ term for $k=n$; $p,q$ and the last row when $\Hspace^k=\{0\}$.

\begin{remark}[Relation to AFW]
Arnold--Falk--Winther write the first row as $\ip{\sigma}{\tau}-\ip{u}{\dd\tau}=-\int_{\Gamma_n}
B_{k-1}(\tau,g_u)$. \code{boundary\_conditions.build\_natural\_terms} returns the boundary
functionals in that (unscaled) convention, i.e.\ $\sigma$-term $=-\sum_{\text{faces}}\int
B_{k-1}(\tau,g_u)$, and \code{EllipticProblem.\_natural\_terms} negates it.
\end{remark}

In block form, with $M_\sigma$ the $V^{k-1}$ mass matrix, $D$ the matrix of
$\ip{\dd\sigma}{v}$, $K$ that of $\ip{\dd u}{\dd v}$ and $C$ that of $\ip{p}{v}$,
\[
  \begin{pmatrix} -M_\sigma & D^T & 0\\ D & K & C\\ 0 & C^T & 0 \end{pmatrix}
  \begin{pmatrix}\sigma\\ u\\ p\end{pmatrix} =
  \begin{pmatrix} b_\sigma\\ b_u\\ 0\end{pmatrix},
\]
which is symmetric and indefinite.

%=================================================================================================
\section{Vector proxies}\label{sec:proxies}
%=================================================================================================

The code represents forms by scalar or vector \emph{proxies}: $0$- and $n$-forms by scalars,
intermediate degrees by vectors in $\R^n$. All proxy maps below are pointwise isometries
(rotations and sign changes), so the $L^2$ inner product of forms is the ordinary
$\int u\cdot v$ (or $\int uv$) of proxies, and the mass matrices are the usual ones.

\subsection{The proxy maps}

\begin{description}
\item[1D] ($\Lambda^0\to\Lambda^1$): $w\mapsto w$, $\ a\,\dd x\mapsto a$.
\item[2D grad-curl] ($\Lambda^1$ proxy in $H(\mathrm{curl})$):
  $a\,\dd x+b\,\dd y\mapsto (a,b)$, $\ c\,\dd x\wedge\dd y\mapsto c$.
\item[2D curl-div] ($\Lambda^1$ proxy in $H(\mathrm{div})$, rotated):
  $a\,\dd x+b\,\dd y\mapsto (-b,\,a)$, $\ c\,\dd x\wedge\dd y\mapsto -c$.
\item[3D] $a_1\dd x+a_2\dd y+a_3\dd z\mapsto(a_1,a_2,a_3)$;
  $\ b_1\,\dd y\wedge\dd z+b_2\,\dd z\wedge\dd x+b_3\,\dd x\wedge\dd y\mapsto(b_1,b_2,b_3)$;
  $\ c\,\dd x\wedge\dd y\wedge\dd z\mapsto c$.
\end{description}

\paragraph{Why the curl-div 2-form proxy carries a minus sign.} In the curl-div complex the
1-form proxy is the rotation $R(a,b)=(-b,a)$ of the Cartesian components. It is forced by
requiring that $\dd$ on 0-forms be $\skewgrad w = (-\partial_y w,\partial_x w)$:
$\dd w = w_x\dd x+w_y\dd y\mapsto(-w_y,w_x)$. With $\dd$ on 1-forms equal to $\div$,
\[
  \dd(a\,\dd x+b\,\dd y) = (\partial_x b-\partial_y a)\,\dd x\wedge\dd y,
  \qquad \div(-b,a) = -(\partial_x b - \partial_y a),
\]
so the scalar proxy of $c\,\dd x\wedge\dd y$ must be $-c$. Flipping the 1-form map instead
would break $\dd w=\skewgrad w$. This sign is invisible in the operators (which are written
directly in proxies) and in the harmonic \emph{spans}; it only fixes the sign of the constant
2-form proxy (Section~\ref{sec:constant}).

\subsection{Operators, adjoints and boundary pairings}

Table~\ref{tab:operators} is implemented by \code{differential\_forms.exterior\_derivative},
\code{codifferential} and \code{boundary\_pairing}. Here $n$ is the outward unit normal,
\[
  \skewgrad w = (-\partial_y w,\ \partial_x w),\qquad
  \curltwo w = \partial_x w_1 - \partial_y w_0,\qquad
  a\times b = a_0 b_1 - a_1 b_0\ \ \text{(2D, scalar)},
\]
matching \code{ufl\_helpers.skewgrad} and \code{ufl\_helpers.curl2D}.

\begin{table}[h]
\centering
\renewcommand{\arraystretch}{1.25}
\begin{tabular}{@{}llll@{}}
\toprule
case ($w\in\Lambda^k$, $z\in\Lambda^{k+1}$) & $\dd w$ & $\cod z$ & $B_k(w,z)$\\
\midrule
1D, $k=0$                         & $\partial_x w$   & $-\partial_x z$                       & $w\,z\,n$\\
2D grad-curl and 3D, $k=0$        & $\grad w$        & $-\div z$                             & $w\,(z\cdot n)$\\
2D curl-div, $k=0$                & $\skewgrad w$    & $-\curltwo z$                         & $w\,(n\times z)$\\
2D grad-curl, $k=1$               & $\curltwo w$     & $-\skewgrad z=(\partial_y z,-\partial_x z)$ & $(n\times w)\,z$\\
3D, $k=1$                         & $\curl w$        & $\curl z$                             & $(n\times w)\cdot z$\\
2D curl-div $k=1$, 3D $k=2$       & $\div w$         & $-\grad z$                            & $(w\cdot n)\,z$\\
any, $k=n$                        & $0$ (\code{None}) & ---                                  & $0$\\
\bottomrule
\end{tabular}
\caption{Proxy operators. $\cod$ in row $k$ is $\cod_{k+1}$, the adjoint of $\dd_k$; it maps
$(k+1)$-form proxies to $k$-form proxies. \code{codifferential(z, j, \dots)} takes the degree
$j$ of its argument $z$ and returns \code{None} for $j=0$.}
\label{tab:operators}
\end{table}

Each row satisfies the proxy version of \eqref{eq:ibp-forms},
\begin{equation}\label{eq:ibp}
  \int_\Omega \dd w\cdot z\,\dd x = \int_\Omega w\cdot \cod z\,\dd x + \int_{\partial\Omega} B_k(w,z)\,\dd s .
\end{equation}
For example, for 2D curl-div $k=0$:
$\int(-w_y z_0 + w_x z_1) = \int w\,(\partial_y z_0-\partial_x z_1) + \int_{\partial\Omega} w\,(n_x z_1-n_y z_0)$,
i.e.\ $\cod z=-\curltwo z$ and $B=w\,(n\times z)$; for 2D grad-curl $k=1$:
$\int(\partial_x w_1-\partial_y w_0)z = \int w\cdot(\partial_y z,-\partial_x z) + \int_{\partial\Omega}(n_x w_1-n_y w_0)z$.
The remaining rows are the classical Green formulas. All signs are checked by
\code{tests/unit/test\_differential\_forms.py::test\_integration\_by\_parts\_identity} (random
smooth fields, all 8 cases with $k<n$), together with $\dd^2=0$ (symbolic and on discrete
functions) and $\cod^2=0$.

\subsection{The Hodge Laplacian of the proxies}

Because the proxy maps are constant isometries on a flat domain, $\Delta_H$ acts on every proxy
as minus the componentwise Laplacian:
\[
\begin{array}{lll}
k=0: & \cod\dd u = -\div\grad u = -\Delta u & (\text{curl-div: } -\curltwo\skewgrad u = -\Delta u),\\[2pt]
\text{2D grad-curl } k=1: & \dd\cod u+\cod\dd u = -\grad\div u + \operatorname{rot}\curltwo u = -\Delta u,\\[2pt]
\text{3D } k=1: & -\grad\div u + \curl\curl u = -\Delta u,\\[2pt]
\text{3D } k=2: & \curl\curl u - \grad\div u = -\Delta u,\\[2pt]
k=n: & \dd\cod u = -\div\grad u \ \ (\text{2D grad-curl: } -\curltwo\skewgrad u) = -\Delta u,
\end{array}
\]
with $\operatorname{rot}c = (\partial_y c,-\partial_x c)$. This is verified for all 12 cases by
\code{test\_hodge\_laplacian\_is\_minus\_vector\_laplacian} and is how
\code{hodge\_laplacian\_expr} turns an exact solution into a source.

\subsection{Discrete spaces}

For the \code{Q} family (tensor-product $Q_r^-$ spaces on quadrilaterals and extruded hexahedra,
order $r$) the discrete complex used by \code{differential\_forms.form\_space} is

\begin{center}
\renewcommand{\arraystretch}{1.2}
\begin{tabular}{@{}lllll@{}}
\toprule
& $V^0$ & $V^1$ & $V^2$ & $V^3$\\
\midrule
1D & CG$_r$ & DG$_{r-1}$ & & \\
2D grad-curl & CG$_r$ & RTCE$_r$ & DG$_{r-1}$ & \\
2D curl-div & CG$_r$ & RTCF$_r$ & DG$_{r-1}$ & \\
3D & CG$_r$ & NCE$_r$ & NCF$_r$ & DG$_{r-1}$\\
\bottomrule
\end{tabular}
\end{center}

Each is a subcomplex ($\dd V^k\subset V^{k+1}$). 3D Q-family meshes are extruded, because NCE/NCF
are only available on extruded hexahedra. RTCE/RTCF/NCE/NCF have no dual basis, so data are put
into them by $L^2$ projection (\code{harmonic\_forms.interpolate\_or\_project}), and inhomogeneous
essential data go through Firedrake's projection fallback in \code{DirichletBC}.

%=================================================================================================
\section{Boundary conditions}\label{sec:bc}
%=================================================================================================

\subsection{Faces}

Every bounded direction contributes two logical faces; periodic directions contribute none:

\begin{center}
\begin{tabular}{@{}lll@{}}
\toprule
direction & low side & high side\\
\midrule
$x$ & \code{left} & \code{right}\\
$y$ & \code{bottom} & \code{top}\\
$z$ & \code{back} & \code{front}\\
\bottomrule
\end{tabular}
\end{center}

In 3D the logical \code{bottom}/\code{top} are the $y$ faces; the $z$ faces are \code{back}/\code{front}
even though on extruded meshes Firedrake calls them \code{"bottom"}/\code{"top"} (measures
\code{ds\_b}/\code{ds\_t}). Firedrake subdomain ids are numbered consecutively $1,2,\dots$ over the
bounded directions (low side first) for intervals, rectangles, \code{PeriodicRectangleMesh} and the
base of extruded meshes; non-extruded 3D meshes (\code{BoxMesh}, \code{PeriodicBoxMesh}) keep the
fixed ids 1--6. \code{boundary\_conditions.verify\_boundary\_registry} checks this geometrically.

\subsection{Essential and natural conditions}

Each face carries one type. With $\tr$ the tangential trace and $\tr\hodge$ the trace of the Hodge
dual (``normal'' trace),
\begin{itemize}
\item \textbf{essential} ($\Gamma_e$): $\tr u = \tr g_u$ and $\tr\sigma=\tr g_\sigma$, imposed strongly by
  \code{DirichletBC} on the natural trace of each space (value for CG, tangential trace for
  $H(\mathrm{curl})$, normal trace for $H(\mathrm{div})$, nothing for DG);
\item \textbf{natural} ($\Gamma_n$): $\tr\hodge u=\tr\hodge g_u$ and $\tr\hodge\dd u = \tr\hodge g_{\dd u}$,
  imposed weakly by the boundary functionals of \eqref{eq:mixed}.
\end{itemize}
The data are passed as \emph{full fields} ($g_u$, $g_\sigma=\cod g_u$, $g_{\dd u}=\dd g_u$, by default
from the exact solution); the traces and pairings are applied by the code. Their meaning per case:

\begin{center}
\renewcommand{\arraystretch}{1.2}
\begin{tabular}{@{}lll@{}}
\toprule
case & essential & natural\\
\midrule
$k=0$ & $u=g$ (Dirichlet) & $\nabla u\cdot n = g$ (Neumann)\\
$k=n$, $\sigma\in H(\mathrm{div})$ or CG (1D) & $\sigma\cdot n=g$ (flux) & $u=g$ (weak Dirichlet)\\
$k=n=2$, grad-curl ($\sigma\in H(\mathrm{curl})$) & $n\times\sigma=g$ & $u=g$ (weak)\\
2D grad-curl $k=1$ ($u\in H(\mathrm{curl})$) & $n\times u$, $\sigma=-\div u$ & $u\cdot n$, $\curltwo u$\\
2D curl-div $k=1$ ($u\in H(\mathrm{div})$) & $u\cdot n$, $\sigma=-\curltwo u$ & $n\times u$, $\div u$\\
3D $k=1$ ($u\in H(\mathrm{curl})$) & $n\times u$, $\sigma=-\div u$ & $u\cdot n$, $n\times\curl u$\\
3D $k=2$ ($u\in H(\mathrm{div})$) & $u\cdot n$, $n\times\sigma$ ($\sigma=\curl u$) & $n\times u$, $\div u$\\
\bottomrule
\end{tabular}
\end{center}

%=================================================================================================
\section{Harmonic forms}\label{sec:harmonic}
%=================================================================================================

\subsection{Definition and the K\"unneth rule}

$\Hspace^k = \{h\in\Lambda^k: \dd h=0,\ \cod h=0,\ \tr h=0 \text{ on }\Gamma_e,\ \tr\hodge h=0\text{ on }\Gamma_n\}$.
For a constant form $\dd x^I$ on a face normal to $x_i$,
\[
  \tr\,\dd x^I = 0 \iff i\in I, \qquad \tr\hodge\dd x^I = 0 \iff i\notin I .
\]
Let $P$ be the periodic directions, $E$ the bounded directions with both faces essential and $N$
those with both faces natural. Then (K\"unneth):
\begin{itemize}
\item if some bounded direction has different types on its two faces, $\Hspace^k=\{0\}$;
\item otherwise $\Hspace^k=\operatorname{span}\{\dd x^I: E\subseteq I\subseteq E\cup P,\ |I|=k\}$,
  of dimension $\binom{|P|}{k-|E|}$.
\end{itemize}
Examples: $k=0$ all-natural gives the constants; $k=n$ all-essential gives the volume form; fully
periodic gives all $\binom{n}{k}$ constant forms; a periodic-$x$ cylinder with natural $y$ faces has
$\Hspace^1=\operatorname{span}\{\dd x\}$, with essential $y$ faces $\Hspace^1=\operatorname{span}\{\dd y\}$.
This is \code{harmonic\_forms.harmonic\_index\_sets}. Constant forms lie in every discrete space
and are discretely harmonic: $\dd h=0$ and $\ip{h}{\dd\tau}=0$ for all $\tau\in V^{k-1}$ with the
essential conditions (checked by \code{test\_harmonic\_basis\_is\_discretely\_harmonic}).

\subsection{Proxies of constant forms}\label{sec:constant}

\code{constant\_form\_values(I, k, n, complex)}, with $I$ in increasing order:

\begin{center}
\renewcommand{\arraystretch}{1.2}
\begin{tabular}{@{}ll@{}}
\toprule
case & proxy of $\dd x^I$\\
\midrule
$k=0$ & $1$\\
$k=n$ & $1$, except 2D curl-div: $\dd x\wedge\dd y\mapsto -1$\\
$H(\mathrm{curl})$ 1-forms (2D grad-curl, 3D) & $\dd x_i\mapsto e_i$\\
2D curl-div 1-forms & $\dd x\mapsto(0,1)$, $\dd y\mapsto(-1,0)$\\
3D 2-forms & $\dd y\wedge\dd z\mapsto e_x$, $\dd x\wedge\dd z\mapsto -e_y$, $\dd x\wedge\dd y\mapsto e_z$\\
\bottomrule
\end{tabular}
\end{center}

The 3D entry $\dd x\wedge\dd z\mapsto -e_y$ is the increasing-order statement of
$\dd z\wedge\dd x\mapsto e_y$. All values are pinned by
$\dd(x_{i_1}\,\dd x^{I\setminus\{i_1\}}) = \dd x^I$ evaluated with the operators of
Table~\ref{tab:operators} (\code{test\_constant\_form\_proxy\_matches\_d\_of\_coordinates}); e.g.\
$\curl(0,0,x)=(0,-1,0)$ and, in curl-div, $\div(x\,(-1,0))=-1$.

\subsection{Orthonormal basis}

\code{orthonormalize\_harmonic\_basis} puts each constant proxy into $V^k$ (interpolation, or $L^2$
projection with LU, which is exact since $V^k$ contains constants) and applies modified
Gram--Schmidt twice in the $L^2$ inner product. All later formulas use this $L^2$-orthonormal basis
$\{h_i\}_{i=1}^m$.

\subsection{Compatibility of exact data}

\begin{proposition}
Let $u$ be smooth, $\sigma=\cod u$, $f=\Delta_H u$, and take exact boundary data. Then the solution
of \eqref{eq:mixed} has $p=0$ (up to discretization of the data).
\end{proposition}
\begin{proof}[Sketch]
Let $h\in\Hspace^k$, so $\dd h=0$ and $\cod h=0$. For the exact fields, \eqref{eq:ibp} applied to
$\ip{\dd\sigma}{h}$ and to $\ip{\dd u}{\dd h}=0$ gives
$\ip{f}{h}=\ip{\dd\sigma+\cod\dd u}{h}=\int_{\partial\Omega}B_{k-1}(\sigma,h)-\int_{\partial\Omega}B_k(h,\dd u)$.
Testing the second row of \eqref{eq:mixed} with $v=h$ and using
$\ip{\dd\sigma_h}{h}=\int_{\partial\Omega}B_{k-1}(\sigma_h,h)$ gives
\[
  \ip{p}{h} = \int_{\partial\Omega} B_{k-1}(\sigma-\sigma_h,h) - \int_{\Gamma_e}B_k(h,\dd u)
            + \int_{\Gamma_n}B_k(h,g_{\dd u}-\dd u) .
\]
On $\Gamma_e$, $\tr\sigma_h=\tr g_\sigma=\tr\sigma$ and $\tr h=0$, so both $\Gamma_e$ terms vanish; on
$\Gamma_n$, $\tr\hodge h=0$ kills $B_{k-1}(\sigma-\sigma_h,h)$ and $g_{\dd u}=\dd u$ kills the last term.
\end{proof}
Hence manufactured problems only need $u\perp\Hspace^k$ (zero mean of every proxy component) for
$u_h$ to approximate $u$ directly; \code{SineProblem} uses full-period factors
$\sin(2\pi m_i s_i+\phi)$ in normalized coordinates $s_i\in[0,1]$, which are periodic and
zero-mean in every direction. Errors are always measured against $u-P_{\Hspace}u$.
\code{IncompatibleSourceProblem} adds $\sum_I a_I\,\dd x^I$ (harmonic $I$ only) to $f$, so that
$p=\sum_I a_I\,\dd x^I$.

%=================================================================================================
\section{Enforcing the harmonic constraint}\label{sec:constraint}
%=================================================================================================

\subsection{\code{lagrange}}

$p=\sum_i p_i h_i$ with one \emph{scalar} $\R$-space block per $h_i$ (Firedrake supports neither
arguments on vector-valued $\R$ spaces nor monolithic matrices with $\R$ blocks). The coupling is
\[
  \sum_i \bigl[p_i\,\ip{h_i}{v} + q_i\,\ip{u}{h_i}\bigr]
\]
(\code{harmonic\_constraint\_form}). The bordered system is solved matrix-free with
right-preconditioned GMRES, preconditioned by the assembled
\[
  a_P = a_0 + \varepsilon\ip{u}{v} + \sum_i p_i q_i ,
\]
where $a_0$ is \eqref{eq:mixed} without the coupling and $\varepsilon=$
\code{lagrange\_regularization} (default 1): an additive fieldsplit with \code{AssembledPC} + LU
on the non-$\R$ blocks and the identity on the $\R$ blocks. On non-harmonic modes the preconditioned
operator has eigenvalues $\lambda/(\lambda+\varepsilon)$ ($\lambda$ the Hodge Laplacian
eigenvalues), and the $m$ harmonic directions coupled to $p$ add $2m$ eigenvalues
$\pm(\varepsilon|\Omega|)^{-1/2}$, so few iterations are needed; $\varepsilon$ should be small
compared to $\pi^2/L^2$.

\subsection{\code{nullspace}}

Drop $p$ and the last row. The operator is singular with kernel
$\{(\sigma,u)=(0,h):h\in\Hspace^k\}$: indeed $\ip{h}{\dd\tau}=0$ and $\dd h=0$. Since the operator is
symmetric (Section~\ref{sec:weak}), the kernel is also the cokernel. The solve proceeds as follows
(\code{EllipticProblem.\_solve\_nullspace}):
\begin{enumerate}
\item \textbf{Lift} the essential data explicitly: $w=w_g+w_0$, where $w_g$ carries the
  \code{DirichletBC} values, and assemble $b=L-a(w_g,\cdot)$, zeroed on the constrained dofs.
\item \textbf{Project} $b$ orthogonal to the harmonic forms: $b\leftarrow b-\sum_i b(h_i)\,Mh_i$ on
  the $u$ block, with $M$ the $V^k$ mass matrix and $Mh_i$ zeroed on the constrained rows. Since
  $\ip{h_i}{h_j}=\delta_{ij}$ this gives $b(h_i)=0$, and $p_i=b(h_i)$ are the reported coefficients
  (\code{project\_rhs\_orthogonal\_to\_harmonic}).
\item \textbf{Solve} for $w_0$ with homogeneous BCs and the nullspace attached
  (\code{VectorSpaceBasis}, $\ell^2$-orthonormalized as PETSc requires), with right-preconditioned
  GMRES preconditioned by the LU of the regularized, non-singular $a+\varepsilon\ip{u}{v}$, or with
  MINRES and the Riesz-map preconditioner $\ip{\sigma}{\tau}+\ip{\dd\sigma}{\dd\tau}+\ip{u}{v}+\ip{\dd u}{\dd v}$.
\item \textbf{Remove} the harmonic part: $u\leftarrow u-\sum_i\ip{u}{h_i}h_i$.
\end{enumerate}
The lifting must come \emph{before} the projection. Firedrake's own lifting would run after it and,
with nonzero $g_\sigma$, reintroduce a harmonic component, since $\ip{h}{\dd\tau_j}\neq0$ for basis
functions $\tau_j$ with nonzero trace on $\Gamma_e$.

\subsection{Consistency of the two approaches}

Both give the same $(\sigma_h,u_h,p_h)$. Testing the \code{lagrange} system with $(0,h_j,0)$ gives
$p_j = b_u(h_j)-\ip{\dd\sigma_h}{h_j}$, and $\ip{\dd(\sigma_h-\sigma_g)}{h_j}=0$ because
$\sigma_h-\sigma_g$ satisfies the homogeneous essential conditions, so $p_j$ equals the coefficient of
the lifted rhs. In practice this requires \emph{the same quadrature in both paths}: the code always uses
an explicit degree (\code{elliptic.quadrature\_degree}, default $2r+2$). With UFL's per-form estimate,
the combined mixed residual of \code{lagrange} and the separately assembled rhs of \code{nullspace}
integrate the source with different rules, and the solutions differ at the quadrature-error level
(observed in 1D, $\sim10^{-5}$). In both approaches $p$ is reported as a Function in $V^k$.

%=================================================================================================
\section{Map from formulas to code}\label{sec:code}
%=================================================================================================

\begin{center}
\renewcommand{\arraystretch}{1.15}
\begin{tabular}{@{}>{\raggedright}p{0.42\textwidth}>{\raggedright\arraybackslash}p{0.52\textwidth}@{}}
\toprule
item & implementation\\
\midrule
spaces $V^k$ & \code{differential\_forms.form\_space} (wraps \code{variables.get\_diff\_form\_space})\\
$\dd$, $\cod$, $B_k$ (Table~\ref{tab:operators}) & \code{exterior\_derivative}, \code{codifferential}, \code{boundary\_pairing}\\
$\Delta_H$ & \code{hodge\_laplacian\_expr}\\
constant-form proxies & \code{constant\_form\_values}, \code{constant\_form\_proxy}, \code{normalize\_index\_set}\\
faces, ids, measures & \code{boundary\_conditions.build\_face\_registry}, \code{measure\_for}\\
essential / natural BCs & \code{build\_essential\_bcs}, \code{build\_natural\_terms}\\
harmonic rule and basis & \code{harmonic\_forms.harmonic\_index\_sets}, \code{harmonic\_basis}, \code{orthonormalize\_harmonic\_basis}\\
\code{lagrange} pieces & \code{harmonic\_multiplier\_space}, \code{harmonic\_constraint\_form}, \code{lagrange\_preconditioner\_form}, \code{lagrange\_solver\_parameters}\\
\code{nullspace} pieces & \code{project\_rhs\_orthogonal\_to\_harmonic}, \code{build\_harmonic\_nullspace}, \code{remove\_harmonic\_component}\\
problem data, MMS & \code{elliptic\_problems} (\code{SineProblem}, \code{PolynomialProblem}, \dots, \code{collect\_boundary\_data})\\
mixed system \eqref{eq:mixed}, solve, errors & \code{elliptic.EllipticProblem}, \code{EllipticConfig}, \code{elliptic\_solver\_parameters}\\
\bottomrule
\end{tabular}
\end{center}

\section*{References}
\begin{itemize}
\item D.~N.~Arnold, R.~S.~Falk, R.~Winther, \emph{Finite element exterior calculus, homological
  techniques, and applications}, Acta Numerica 15 (2006) 1--155.
\item D.~N.~Arnold, R.~S.~Falk, R.~Winther, \emph{Finite element exterior calculus: from Hodge theory
  to numerical stability}, Bull.\ AMS 47 (2010) 281--354.
\item D.~N.~Arnold, \emph{Finite Element Exterior Calculus}, CBMS-NSF Regional Conference Series in
  Applied Mathematics 93, SIAM (2018).
\end{itemize}

\end{document}
